% =============================================================================
% EIT-EMP-1: EMPIRICAL GAMMA EMERGENCE SPECIFICATION
% =============================================================================
%
% Erweiterung zu Appendix IE (CORE-EIT), Axiom EIT-15
% Zur Integration in Appendix IE oder als eigenständiges Supplement
%
% Version: 1.0 (Januar 2026)
% =============================================================================

\section{EIT-EMP-1: Empirical Gamma Emergence}
\label{sec:ie-emp1-gamma}

\begin{tcolorbox}[colback=purple!5!white, colframe=purple!75!black,
    title=\textbf{Empirical Extension EIT-EMP-1: Gamma Emergence Specification}]

\textbf{Extends:} Axiom EIT-15 (General Complementarity Function)

\textbf{Purpose:} EIT-15 defines that $\gamma: [0,1]^9 \times [0,1]^9 \rightarrow [-1,1]$ exists.
EIT-EMP-1 specifies \textit{how} $\gamma$ is concretely computed through:
\begin{enumerate}[nosep]
    \item Literature-derived baseline values $\gamma_0$
    \item Context-sensitivity parameters $\beta$ that emerge via LLMMC
    \item A continuous update mechanism (feedback loop)
\end{enumerate}

\textbf{Status:} Empirical extension (data-dependent)

\end{tcolorbox}

%------------------------------------------------------------------------------
\subsection{The 7 Emergence Principles}
\label{subsec:emp1-principles}
%------------------------------------------------------------------------------

The emergent intervention system follows seven principles derived from EIT:

\begin{table}[htbp]
\centering
\caption{The 7 Emergence Principles (EIT-EMP-1)}
\label{tab:emergence-principles}
\renewcommand{\arraystretch}{1.3}
\begin{tabular}{@{}clll@{}}
\toprule
\textbf{\#} & \textbf{Principle} & \textbf{Source} & \textbf{What Emerges?} \\
\midrule
E1 & Continuous Space & EIT-1 & No types, only vectors in $[0,1]^k$ \\
E2 & Context Transformation & EIT-2 & Effect $\Delta S = f(I, S_0, \Psi)$ \\
E3 & Cluster Emergence & EIT-3 & ``Types'' as density regions \\
E4 & Journey Emergence & EIT-4 & Phase $\varphi = h(S_0)$ \\
E5 & Horizon Emergence & EIT-5 & Effect duration $\tau = g(\vec{I})$ \\
E6 & Gamma Emergence & EIT-15 + EMP-1 & $\gamma_{ij}(\Psi) = \gamma_0 + \sum \beta \cdot \Psi$ \\
E7 & Mode Emergence & EIT-10 & E-Full/Partial/None from $D_{suff}$ \\
\bottomrule
\end{tabular}
\end{table}

%------------------------------------------------------------------------------
\subsection{The Three-Layer Gamma Architecture}
\label{subsec:emp1-architecture}
%------------------------------------------------------------------------------

\begin{tcolorbox}[colback=blue!5!white, colframe=blue!75!black,
    title=\textbf{Definition EIT-EMP-1-D1: Three-Layer Gamma Architecture}]

The complementarity function $\gamma$ is computed through three layers:

\textbf{Layer 1: Literature Baseline ($\gamma_0$)}
\begin{itemize}[nosep]
    \item Source: Meta-analyses, experimental studies
    \item Status: Fixed after initial calibration
    \item Update: Annually with new meta-evidence
\end{itemize}

\textbf{Layer 2: Context Sensitivity ($\beta$)}
\begin{itemize}[nosep]
    \item Source: LLMMC from observed interventions
    \item Status: Emerges continuously
    \item Update: Bayesian update with each new observation
\end{itemize}

\textbf{Layer 3: Emergent Computation}
\begin{equation}
\boxed{\gamma_{ij}(\Psi) = \gamma_{0,ij} + \sum_{k=1}^{8} \beta_{ijk} \cdot \Psi_k}
\label{eq:gamma-emergence}
\end{equation}

where $\Psi_k$ are the 8 context dimensions (economic, social, temporal, spatial, institutional, cultural, technological, environmental).

\end{tcolorbox}

%------------------------------------------------------------------------------
\subsection{Baseline Parameters from Literature}
\label{subsec:emp1-baseline}
%------------------------------------------------------------------------------

\begin{tcolorbox}[colback=gray!5!white, colframe=gray!75!black,
    title=\textbf{EIT-EMP-1-D2: The 48 Baseline Parameters}]

The baseline $\gamma_0$ values are organized in four categories:

\textbf{1. Referent Complementarity (1 parameter):}
\begin{itemize}[nosep]
    \item $\gamma_0(\text{Self}, \text{Comm}) = -0.30$ \quad (Egocentric Bias, AWX-A34)
\end{itemize}

\textbf{2. Utility-Type Complementarity (3 parameters):}
\begin{itemize}[nosep]
    \item $\gamma_0(\text{INU}, \text{KNU}) = +0.25$ \quad (AWX-A38)
    \item $\gamma_0(\text{INU}, \text{IDN}) = +0.40$ \quad (AWX-A39)
    \item $\gamma_0(\text{KNU}, \text{IDN}) = +0.35$ \quad (AWX-A40)
\end{itemize}

\textbf{3. Domain Complementarity (15 parameters):}
\begin{center}
\small
\begin{tabular}{@{}lcl@{}}
$\gamma_0(F, E) = +0.15$ & $\gamma_0(E, P) = +0.35$ & $\gamma_0(P, D) = +0.15$ \\
$\gamma_0(F, P) = +0.20$ & $\gamma_0(E, S) = +0.40$ & $\gamma_0(P, Ex) = +0.10$ \\
$\gamma_0(F, S) = \mathbf{-0.25}$ & $\gamma_0(E, D) = +0.25$ & $\gamma_0(S, D) = +0.35$ \\
$\gamma_0(F, D) = +0.30$ & $\gamma_0(E, Ex) = +0.45$ & $\gamma_0(S, Ex) = +0.50$ \\
$\gamma_0(F, Ex) = +0.10$ & $\gamma_0(P, S) = +0.20$ & $\gamma_0(D, Ex) = +0.40$ \\
\end{tabular}
\end{center}

\textbf{Note:} $\gamma_0(F, S) = -0.25$ is the \textbf{Crowding-Out} baseline (sign fixed by Theorem EIT-T9).

\textbf{4. Valence Complementarity (1 parameter):}
\begin{itemize}[nosep]
    \item $\gamma_0(\text{pos}, \text{neg}) = -0.50$ \quad (Loss Aversion, $\lambda \approx 2$)
\end{itemize}

\textbf{5. Block-Level 10C (21 parameters) and Cross-Block (8 parameters):} See YAML specification.

\textbf{Total: 48 baseline parameters}

\end{tcolorbox}

%------------------------------------------------------------------------------
\subsection{Constraint: Sign Preservation (EIT-T9)}
\label{subsec:emp1-constraint}
%------------------------------------------------------------------------------

\begin{tcolorbox}[colback=red!5!white, colframe=red!75!black,
    title=\textbf{EIT-EMP-1-C1: Sign Preservation Constraint}]

Per Theorem EIT-T9, the \textbf{sign} of crowding-out is context-invariant:

\begin{equation}
\gamma(F, S \;|\; \Psi) < 0 \quad \forall \Psi \in [0,1]^8
\end{equation}

\textbf{Enforcement:} The sensitivity parameters $\beta$ are constrained such that:
\begin{equation}
\gamma_0(F, S) + \sum_{k=1}^{8} \beta_{FS,k} \cdot \Psi_k < 0 \quad \forall \Psi \in [0,1]^8
\end{equation}

\textbf{Interpretation:} Context can \textit{weaken} crowding-out (make $\gamma$ closer to 0) but never \textit{reverse} it (make $\gamma > 0$).

\end{tcolorbox}

%------------------------------------------------------------------------------
\subsection{The 80-Dimensional Awareness Space}
\label{subsec:emp1-awareness}
%------------------------------------------------------------------------------

For AWARE-focused interventions, the space expands to 80 dimensions:

\begin{tcolorbox}[colback=blue!5!white, colframe=blue!75!black,
    title=\textbf{EIT-EMP-1-D3: The 80D Awareness Intervention Space}]

\textbf{Block A: Other 10C Dimensions (8D)}
\begin{equation}
D_1 = \text{WHO}, \; D_2 = \text{WHAT}, \; D_3 = \text{HOW}, \; D_4 = \text{WHEN}, \; D_5 = \text{WHERE}, \; D_6 = \text{READY}, \; D_7 = \text{STAGE}, \; D_8 = \text{HIER}
\end{equation}

\textbf{Block B: AWARE Detail (72D)}
\begin{equation}
72 = 2 \times 3 \times 6 \times 2 = \underbrace{\text{Self/Comm}}_{\text{Referent}} \times \underbrace{\text{INU/KNU/IDN}}_{\text{Utility}} \times \underbrace{\text{FEPSDE}}_{\text{Domain}} \times \underbrace{\text{pos/neg}}_{\text{Valence}}
\end{equation}

\textbf{Enumeration:}
\begin{align*}
D_9 &= \text{Self\_INU\_F\_pos} \\
D_{10} &= \text{Self\_INU\_F\_neg} \\
&\vdots \\
D_{80} &= \text{Comm\_IDN\_Ex\_neg}
\end{align*}

\textbf{Total: 8 + 72 = 80 dimensions}

\end{tcolorbox}

%------------------------------------------------------------------------------
\subsection{Gamma Computation for Detail Dimensions}
\label{subsec:emp1-computation}
%------------------------------------------------------------------------------

\begin{tcolorbox}[colback=green!5!white, colframe=green!75!black,
    title=\textbf{EIT-EMP-1-A1: Detail Gamma Computation}]

For two AWARE-detail dimensions $D_i$ and $D_j$, the complementarity is computed from factor-level $\gamma_0$ values:

\begin{equation}
\gamma(D_i, D_j) = \frac{1}{4} \left[ \gamma_{\text{ref}}(r_i, r_j) + \gamma_{\text{util}}(u_i, u_j) + \gamma_{\text{dom}}(d_i, d_j) + \gamma_{\text{val}}(v_i, v_j) \right]
\end{equation}

where $(r, u, d, v)$ are the factor indices (referent, utility-type, domain, valence) of each dimension.

\textbf{Example:}
\begin{align*}
\gamma(\text{Self\_INU\_F\_pos}, \text{Comm\_KNU\_S\_neg}) &= \frac{1}{4} \left[ \gamma(\text{Self}, \text{Comm}) + \gamma(\text{INU}, \text{KNU}) + \gamma(F, S) + \gamma(\text{pos}, \text{neg}) \right] \\
&= \frac{1}{4} \left[ -0.30 + 0.25 + (-0.25) + (-0.50) \right] \\
&= \frac{1}{4} \times (-0.80) = -0.20
\end{align*}

\end{tcolorbox}

%------------------------------------------------------------------------------
\subsection{LLMMC Protocol for Beta Estimation}
\label{subsec:emp1-llmmc}
%------------------------------------------------------------------------------

\begin{tcolorbox}[colback=yellow!5!white, colframe=yellow!75!black,
    title=\textbf{EIT-EMP-1-P1: LLMMC Protocol for Sensitivity Estimation}]

\textbf{Input:}
\begin{itemize}[nosep]
    \item Intervention description $\mathcal{I}$ with known profile $\vec{I}$
    \item Context vector $\Psi$ at time of implementation
    \item Observed outcome (effect size, success/failure)
\end{itemize}

\textbf{LLMMC Query:}
\begin{quote}
\textit{``Given intervention $\mathcal{I}$ in context $\Psi$ with outcome $O$: Which $\gamma$ values were likely active? How strong was the context modulation compared to baseline?''}
\end{quote}

\textbf{Output:}
\begin{itemize}[nosep]
    \item Active $\gamma$ pairs with estimated values
    \item Deviation from baseline: $\Delta\gamma = \gamma_{\text{observed}} - \gamma_0$
    \item Implied $\beta$ contribution: $\beta_k \approx \Delta\gamma / \Psi_k$
    \item Confidence level
\end{itemize}

\textbf{Update Rule (Bayesian):}
\begin{equation}
\beta_k^{(n+1)} = \beta_k^{(n)} + \alpha \cdot (\beta_k^{\text{implied}} - \beta_k^{(n)})
\end{equation}

where $\alpha$ is a learning rate weighted by confidence.

\end{tcolorbox}

%------------------------------------------------------------------------------
\subsection{Summary: What is Fixed vs. What Emerges}
\label{subsec:emp1-summary}
%------------------------------------------------------------------------------

\begin{table}[htbp]
\centering
\caption{Fixed vs. Emergent Components in EIT-EMP-1}
\label{tab:fixed-vs-emergent}
\renewcommand{\arraystretch}{1.3}
\begin{tabular}{@{}lll@{}}
\toprule
\textbf{Component} & \textbf{Status} & \textbf{Source} \\
\midrule
\multicolumn{3}{l}{\textit{FIXED (Input)}} \\
\midrule
10C Dimensions & Fixed & Framework definition \\
$\gamma_0$ Baseline (48 values) & Fixed & Literature \\
Sign of $\gamma(F,S)$ & Fixed (negative) & Theorem EIT-T9 \\
Emergence principles (E1-E7) & Fixed & EIT Axioms \\
\midrule
\multicolumn{3}{l}{\textit{EMERGENT (Output)}} \\
\midrule
$\beta$ Sensitivities & Emerges & LLMMC from observations \\
$\gamma(\Psi)$ Context-specific & Emerges & Computation from $\gamma_0 + \beta \cdot \Psi$ \\
Cluster structure & Emerges & Density in observation space \\
Number of clusters $K$ & Emerges & Data-dependent \\
\bottomrule
\end{tabular}
\end{table}

%------------------------------------------------------------------------------
\subsection{Integration with Existing EIT}
\label{subsec:emp1-integration}
%------------------------------------------------------------------------------

\begin{tcolorbox}[colback=purple!5!white, colframe=purple!75!black,
    title=\textbf{Relationship to Existing EIT Components}]

\textbf{EIT-15 (General Complementarity Function):}
\begin{itemize}[nosep]
    \item EIT-15 defines: $\gamma: [0,1]^9 \times [0,1]^9 \rightarrow [-1,1]$ exists
    \item EIT-EMP-1 provides: The concrete parametrization $\gamma(\Psi) = \gamma_0 + \sum \beta \cdot \Psi$
\end{itemize}

\textbf{EIT-T9 (Crowding-Out Universality):}
\begin{itemize}[nosep]
    \item EIT-T9 proves: $\gamma(F,S) < 0$ for all contexts
    \item EIT-EMP-1 respects: Constraint on $\beta$ ensures sign preservation
\end{itemize}

\textbf{EIT-10 (Data-Mode Correspondence):}
\begin{itemize}[nosep]
    \item EIT-10 defines: Mode selection based on $D_{\text{suff}}$
    \item EIT-EMP-1 extends: $\beta$ estimation quality affects $D_{\text{suff}}$
\end{itemize}

\end{tcolorbox}

%------------------------------------------------------------------------------
% END OF EIT-EMP-1
%------------------------------------------------------------------------------
